\section{\texorpdfstring{$f$}{f}-divergence occupancy matching: \texorpdfstring{$f$}{f}-PG with a fixed expert target}\label{sec:m-fpg}

FPG minimizes $D_f(p_\theta\Vert p_g)$ between the agent's state visitation and a Dirac goal target $p_g=\delta_g$ (FPG Eq.~4, p.~4).
For imitation we replace the goal by a \emph{fixed expert distribution} and the visitation by the $H$-step next-state occupancy.
This is an adaptation, not a reproduction of FPG's experiments: the gradient below is derived for this target, and FPG's Dirac-target optimality theorems do not transfer (\cref{rem:m-dirac}).
Complete proofs are in \cref{sec:div,sec:fpg}.

\subsection{Divergences: the facts used}

\begin{proposition}[nonnegativity, orientation and smoothing; \cref{prop:div-nonneg,prop:div-generators,prop:div-perspective,prop:div-smooth}]\claim{DIV-1--DIV-3}{I+D}\label{prop:m-div}
Let $f$ be convex with $f(1)=0$.
\begin{enumerate}[nosep,label=(\alph*)]
\item $D_f(p\Vert q)\ge0$, with equality if and only if $p=q$ when $f$ is strictly convex at $1$; this holds for the three generators of \cref{sec:notation}. $0\le\JS\le\log2$; $\KL(p\Vert q)<\infty$ if and only if $\supp p\subseteq\supp q$.
\item (orientation) With the perspective $f^\circ(u):=uf(1/u)$, $D_f(p\Vert q)=D_{f^\circ}(q\Vert p)$, including when $p$ or $q$ has zeros. $(u\log u)^\circ=-\log u$ and $f_{\JS}^\circ=f_{\JS}$.
\item (smoothing) For a strictly positive reference law $\upsilon$ and $\varepsilon\in(0,1]$, every ratio $p_\varepsilon/q_\varepsilon$ lies in $[\varepsilon\upsilon_{\min},1/(\varepsilon\upsilon_{\min})]$, so $D_f(p_\varepsilon\Vert q_\varepsilon)$ is finite. If $p=d_\phi$ depends smoothly on $\phi$ and $q$ is fixed, $\nabla_\phi D_f((d_\phi)_\varepsilon\Vert q_\varepsilon)=(1-\varepsilon)\sum_xf'\big((d_\phi)_\varepsilon(x)/q_\varepsilon(x)\big)\nabla_\phi d_\phi(x)$.
\end{enumerate}
\end{proposition}

Part (b) translates FIRL's expert-first objective into our learner-first convention: FIRL's ``FKL'' ($f=u\log u$, expert first) is our $\KL(q\Vert d)$, and its ``RKL'' is our $\KL(d\Vert q)$ (\cref{cor:div-dictionary}).
Smoothing changes the objective: the minimizers of the smoothed divergence need not be those of the unsmoothed one, and the evaluation always reports the unsmoothed divergence (\cref{rem:div-smooth}).

\subsection{The objective and the exact gradient}

Let $\{B_1,\dots,B_n\}$ be a finite partition of $\cS$ into \emph{cells}: the states of a finite maze, or the \vn{env.hist_cells}{356} histogram cells of positive free area of the continuous maze (\cref{prop:env-cells}).
For a parametric policy $\pi_\phi$ write $d_\phi(b):=d^+_{\pi_\phi,H}(B_b)$, and let $q$ be a \emph{fixed} probability vector on the cells.
The objective is
\begin{equation}\label{eq:m-fpg-J}
J(\phi):=D_f(d_\phi\Vert q)=\sum_bq(b)\,f\big(d_\phi(b)/q(b)\big).
\end{equation}
We assume (\cref{ass:fpg}): $\pi_\phi(a\mid s)>0$, $\phi\mapsto\pi_\phi(a\mid s)$ is $C^1$, and the score $\chi_k:=\nabla_\phi\log\pi_\phi(a_k\mid s_k)$ is bounded locally uniformly in $\phi$; $f$ is convex and $C^1$ on $(0,\infty)$; $q(b)>0$ and $d_\phi(b)>0$ on every cell, or both arguments are smoothed.
Tabular softmax policies satisfy the score condition with $\lVert\chi_k\rVert_2\le\sqrt2$ (\cref{lem:fpg-softmax}); MLP policies satisfy it with $C^1$ activations on a bounded state space (\cref{rem:fpg-mlp}).
Define the \emph{$f$-PG signal} and the \emph{signal-to-go}
\[
g_\phi(b):=f'\big(d_\phi(b)/q(b)\big),\quad g_\phi(s):=g_\phi(b(s)),\qquad G_k:=\sum_{t=k}^{H-1}\gamma^{t-k}g_\phi(s_{t+1}).
\]

\begin{theorem}[exact gradient, causal form; \cref{thm:fpg-grad,thm:fpg-causal}]\claim{FPG-1, FPG-2}{D}\label{thm:m-fpg}
Under these assumptions $J$ is $C^1$ and
\begin{align}
\nabla d_\phi(b)&=Z_H^{-1}\sum_{t=0}^{H-1}\gamma^t\,\E_\phi\Big[\ind[s_{t+1}\in B_b]\sum_{k=0}^{t}\chi_k\Big],\label{eq:m-fpg-dd}\\
\nabla J(\phi)&=\sum_bg_\phi(b)\,\nabla d_\phi(b)=Z_H^{-1}\,\E_\phi\Big[\sum_{k=0}^{H-1}\gamma^k\,\chi_k\,G_k\Big].\label{eq:m-fpg-grad}
\end{align}
\end{theorem}

\begin{proof}
\emph{Step 1: likelihood ratio.} Let $h$ be a bounded measurable function of $(s_0,a_0,\dots,s_H)$ that does not depend on $\phi$. Writing $\E_\phi h$ as an integral of $h\prod_{k<H}\pi_\phi(a_k\mid s_k)$ against the $\phi$-free measure $\mu_0\otimes\prod_kP(\dd s_{k+1}\mid s_k,a_k)$ (times counting measure on actions), the integrand is $C^1$ in $\phi$ with gradient $h\prod_k\pi_\phi(a_k\mid s_k)\sum_k\chi_k$, bounded by $H C\lVert h\rVert_\infty$ on a ball around $\phi$. The mean value theorem and dominated convergence allow differentiation under the integral:
\[
\nabla_\phi\E_\phi h=\E_\phi\Big[h\sum_{k=0}^{H-1}\chi_k\Big].
\]
\emph{Step 2: causality.} Let $\cF_k$ be generated by $(s_0,a_0,\dots,s_k)$. Given $\cF_k$, $a_k\sim\pi_\phi(\cdot\mid s_k)$, so $\E[\chi_k\mid\cF_k]=\sum_a\pi_\phi(a\mid s_k)\nabla\log\pi_\phi(a\mid s_k)=\nabla\sum_a\pi_\phi(a\mid s_k)=0$. Hence $\E[Y\chi_k]=\E[Y\,\E[\chi_k\mid\cF_k]]=0$ for every bounded $\cF_k$-measurable $Y$.
\emph{Step 3: occupancy gradient.} Apply Step~1 to $h=\ind[s_{t+1}\in B_b]$. For $k\ge t+1$ the indicator is $\cF_k$-measurable, so these terms vanish by Step~2; multiplying by $\gamma^t/Z_H$ and summing over $t$ gives \eqref{eq:m-fpg-dd}.
\emph{Step 4: chain rule.} Since $d_\phi(b)/q(b)\in(0,\infty)$, $\nabla[q(b)f(d_\phi(b)/q(b))]=q(b)f'(d_\phi(b)/q(b))\nabla d_\phi(b)/q(b)=g_\phi(b)\nabla d_\phi(b)$; sum over $b$.
\emph{Step 5: reordering.} Substitute \eqref{eq:m-fpg-dd} and use $\sum_bg_\phi(b)\ind[s_{t+1}\in B_b]=g_\phi(s_{t+1})$:
$\nabla J=Z_H^{-1}\sum_{t<H}\gamma^t\E[g_\phi(s_{t+1})\sum_{k\le t}\chi_k]$. The index set $\{0\le k\le t\le H-1\}$ can be summed over $t$ first or over $k$ first, and $\gamma^t=\gamma^k\gamma^{t-k}$; this gives the causal form.
\end{proof}

\paragraph{Signs and weights.}
The update is a \emph{descent}, $\phi\leftarrow\phi-\eta\widehat{\nabla J}$: in REINFORCE terms \citep{williams1992simple,sutton2000policy} the per-step ``reward'' is $-g_\phi(s_{t+1})$. Since $f'$ is nondecreasing, $g_\phi$ is larger on cells that the learner over-visits relative to the target, and descent lowers the probability of actions whose signal-to-go is above the baseline (a first-order heuristic, not a theorem).
The score at step $k$ carries the weight $\gamma^k$ and the signal at $s_{t+1}$ the weight $\gamma^{t-k}$ inside $G_k$. Dropping $\gamma^k$, common in discounted policy gradients, does \emph{not} give $\nabla J$; FPG's Eqs.~(12)--(13) also differ from \eqref{eq:m-fpg-grad}, while its on-policy Eq.~(10), re-indexed, equals it \audit{F5}.
The signal $g_\phi$ is not a reward: it depends on $\phi$ through $d_\phi$ and is recomputed at every iteration.

\begin{corollary}[baselines and Monte Carlo estimator; \cref{prop:fpg-baseline,prop:fpg-mc}]\claim{FPG-3, FPG-4}{I}\label{cor:m-fpg-mc}
(a) For every bounded $\cF_k$-measurable $\beta_k$ (a function of the history up to $s_k$, for example of $(k,s_k)$), $\nabla J=Z_H^{-1}\E_\phi[\sum_k\gamma^k\chi_k(G_k-\beta_k)]$; adding a constant to $g_\phi$ leaves $\nabla J$ unchanged; a ``baseline'' that depends on $a_k$ changes it in general.
(b) With the \emph{exact} signal, the average $\hat\nabla_N$ of this expression over $N$ independent trajectories is unbiased, has covariance $\mathrm{Cov}(\hat\nabla_1)/N$ and converges almost surely. If $\nabla J(\phi)\ne0$, an exact step $\phi-\eta\nabla J(\phi)$ decreases $J$ for small $\eta$.
\end{corollary}
Nothing more is claimed about optimization: $J$ is not convex in $\phi$, and no statement is made about the stationary points that descent may reach.

\subsection{From the exact gradient to the implemented estimator}\label{sec:m-fpg-approx}

The signal $g_\phi$ requires the unknown $d_\phi$; the implementation replaces it by a \emph{plug-in} estimate from sampled trajectories.
The next result says exactly what this costs.

\begin{proposition}[plug-in signals, smoothing and cross-fitting; \cref{prop:fpg-approx,prop:env-crossfit}]\claim{FPG-5, FPG-9}{D}\label{prop:m-fpg-approx}
\begin{enumerate}[nosep,label=(\alph*)]
\item If the signal $\hat g$ is a random function on cells \emph{independent} of the trajectories in the score sums, then $\E[\hat\nabla_N\mid\hat g]=\sum_b\hat g(b)\nabla d_\phi(b)$ and $\lVert\E[\hat\nabla_N\mid\hat g]-\nabla J\rVert\le\max_b\lvert\hat g(b)-g_\phi(b)\rvert\,C\kappa_H$ with $\kappa_H:=Z_H^{-1}\sum_{t<H}(t+1)\gamma^t$.
\item The plug-in signal is biased in general, $\E f'(\hat d/q)\ne f'(d_\phi/q)$, because $f'$ is nonlinear; the resulting gradient bias is $\sum_b(\E\hat g(b)-g_\phi(b))\nabla d_\phi(b)$. It vanishes if the signal bias is constant across cells, because $\sum_b\nabla d_\phi(b)=0$, but constancy is not necessary; in general no unbiasedness of the gradient estimator is claimed. If $\hat g$ is computed from the \emph{same} trajectories as the scores, (a) fails: for a one-step example with $\KL(d\Vert q)$ the same-batch estimator has mean $0$ while $\nabla J=3\log3/16>0$.
\item With both arguments smoothed, \cref{thm:m-fpg} and \cref{cor:m-fpg-mc} hold verbatim for $J_\varepsilon(\phi):=D_f((d_\phi)_\varepsilon\Vert q_\varepsilon)$ with the signal $(1-\varepsilon)f'\big((d_\phi)_\varepsilon/q_\varepsilon\big)$, and no positivity of $d_\phi$ is needed.
\item (cross-fitting) Split the $N$ trajectories of an iteration into folds $\mathcal I_1,\dots,\mathcal I_{n_{\mathrm f}}$, $n_{\mathrm f}\ge2$, and give the trajectories of fold $j$ a signal $\hat g_{-j}$ and a time baseline $\beta_{-j}$ computed from the other folds and the demonstrations only. Conditionally on the other folds, each term of fold $j$ has mean $\Gamma_\phi(\hat g_{-j}):=\nabla_\phi\int\hat g_{-j}\dd d^+_{\pi_\phi,H}$; the baseline contributes nothing; for cell signals $\Gamma_\phi(\hat g_{-j})=\sum_b\hat g_{-j}(b)\nabla d_\phi(b)$ and (a) applies.
\item For a kernel-density signal evaluated at points, no exact cell signal exists, and no bound between $\Gamma_\phi(\hat g_{-j})$ and a gradient of $J$ is claimed (\cref{rem:fpg-kde}).
\end{enumerate}
\end{proposition}

Cross-fitting \citep[in the spirit of][]{chernozhukov2018double} removes the same-batch dependence of (b), not the plug-in bias. Conditionally on the other fold, the sampled update is an unbiased but noisy estimate of the gradient of the surrogate $\int\hat g_{-j}\dd d_\phi$, in which the fitted signal is held fixed. It is not an unbiased estimate of $\nabla J$: the plug-in bias of (b) remains, and for the kernel-density signal no bound relates the surrogate gradient to $\nabla J$ (e).

\begin{algobox}{Algorithm 1: implemented $f$-PG iteration (\texttt{il\_memoir/fpg.py}, \texttt{FPGLearner.step})}
\small
\textbf{Fixed before learning.} The target $\hat q_E$ from the demonstration \emph{states} (weights $\gamma^t/(N_EZ_H)$ on $s^i_{t+1}$), smoothed with $\varepsilon=0.01$: counts on the \vn{env.cells.4r}{40} or \vn{env.cells.open}{49} states (finite), or a cut-and-normalize weighted Gaussian KDE with bandwidth $h=0.05$ and the reference density $1/\Leb(\cS)$ (continuous; \cref{prop:env-est}). The learner never receives $P$, rewards, goals or expert actions.\\[2pt]
\textbf{Policy.} Finite: tabular softmax logits, initialized at $0$ (uniform policy). Continuous: categorical MLP on raw $(x,y)$ with two $\tanh$ hidden layers of width $64$, output layer scaled by $0.01$ at initialization.\\[2pt]
\textbf{Each iteration.}
\begin{enumerate}[nosep,leftmargin=1.3em]
\item Collect $N=20$ trajectories of $H=100$ transitions from $\mu_0$ with $\pi_\phi$; all $2{,}000$ transitions are counted.
\item Assign trajectory $i$ to fold $i\bmod 2$. For each fold $j$, estimate the learner occupancy $\hat d_{-j}$ from the other fold with the same weights and estimator as the target; set $\hat g_{-j}(s^i_{t+1})=(1-\varepsilon)f'\big(\hat d_{-j,\varepsilon}(s^i_{t+1})/\hat q_{E,\varepsilon}(s^i_{t+1})\big)$ for $i\in\mathcal I_j$, the signal-to-go $\hat G^i_k=\sum_{t\ge k}\gamma^{t-k}\hat g_{-j}(s^i_{t+1})$, and the baseline $\beta_{-j}(k)$, the mean signal-to-go of the other fold at step $k$.
\item Minimize the surrogate $L(\phi)=\frac1N\sum_i\sum_{k<H}\frac{\gamma^k}{Z_H}(\hat G^i_k-\beta_{-j(i)}(k))\log\pi_\phi(a^i_k\mid s^i_k)$, whose gradient is the estimator of \cref{prop:m-fpg-approx}(d), by one Adam step \citep{kingma2015adam} with learning rate $0.3$ (finite) or $0.003$ (continuous).
\end{enumerate}
Budget: 25 iterations ($50{,}000$ transitions) in the finite mazes, 50 iterations ($100{,}000$) in the continuous maze. No PPO clipping and no entropy bonus are used.
\end{algobox}

\begin{table}[ht]
\centering\small
\caption{$f$-PG: exact objects and their implemented approximations. \textsf{[I]} implementation approximation, \textsf{[C]} change of objective.}\label{tab:fpg-ledger}
\begin{tabularx}{\textwidth}{@{}>{\raggedright\arraybackslash}p{2.6cm}>{\raggedright\arraybackslash}p{3.4cm}>{\raggedright\arraybackslash}p{3.6cm}>{\raggedright\arraybackslash}X@{}}
\toprule
Quantity & Exact object & Implemented & Consequence\\
\midrule
objective & $D_f(d_\phi\Vert q_E)$ & $D_f((d_\phi)_\varepsilon\Vert\hat q_{E,\varepsilon})$, 4 demonstrations \textsf{[C]} & the gradient formula of \cref{thm:m-fpg} holds for the smoothed objective (\cref{prop:m-fpg-approx}(c)); its minimizers need not be those of the unsmoothed one\\
signal & $f'(d_\phi/q)$ & plug-in from the other fold \textsf{[I]} & biased (b); conditional mean $\Gamma_\phi(\hat g_{-j})$ (d); bound (a) for counts\\
continuous signal & cell signal & KDE at the sampled points \textsf{[I]} & no bound to a gradient of $J$ (e); Euclidean truncation ignores walls (\cref{rem:env-kde})\\
gradient & expectation & 20 trajectories per step \textsf{[I]} & unbiased for $\nabla J$ only with the exact signal (\cref{cor:m-fpg-mc}(b)); with the fitted signal, conditionally unbiased for the surrogate gradient only (d)\\
step & exact small step & Adam, fixed learning rate \textsf{[I]} & no descent or convergence guarantee\\
policy class & any $C^1$ family & tabular softmax; $\tanh$ MLP & a deterministic expert is not representable with finite logits (\cref{prop:m-fpg-opt})\\
\bottomrule
\end{tabularx}
\end{table}

\subsection{State--action occupancies}

\begin{proposition}[state--action extension; \cref{prop:fpg-sa,cor:fpg-sa-finite}]\claim{FPG-6}{D}\label{prop:m-fpg-sa}
In a finite MDP, for $J^{\SA}(\phi):=D_f(d^{\SA}_{\pi_\phi,H}\Vert q^{\SA})$ with $q^{\SA}$ fixed and both positive on the computation set,
\[
\nabla J^{\SA}=Z_H^{-1}\E_\phi\Big[\sum_{k=0}^{H-1}\gamma^k\chi_k\,G^{\SA}_k\Big],\qquad G^{\SA}_k:=\sum_{t=k}^{H-1}\gamma^{t-k}f'\Big(\frac{d^{\SA}_\phi}{q^{\SA}}(s_t,a_t)\Big).
\]
The current-step term $t=k$ is \emph{not} a baseline, because $\ind[a_t=a]$ is not $\cF_t$-measurable.
Against an unsmoothed deterministic expert, $\JS$ and $\KL(q^{\SA}_E\Vert d^{\SA}_\phi)$ are finite while $\KL(d^{\SA}_\phi\Vert q^{\SA}_E)=+\infty$ for every positive policy; smoothing both arguments makes both $\KL$ directions finite.
\end{proposition}

We do not optimize $J^{\SA}$: the learners use states only. The finite evaluation reports $\JS(d^{\SA}_{\pi,H},q^{\SA}_E)$ as a diagnostic of how far state matching is from action matching.

\subsection{The state-entropy reading and the optimum}

\begin{proposition}[state-MaxEnt identity; FPG Lemma~5.1, p.~6; \cref{prop:fpg-maxent}]\claim{FPG-7}{S+D}\label{prop:m-maxent}
If $\supp d\subseteq\supp q$, then $\KL(d\Vert q)=-\Ent(d)-\E_d[\log q]$. Minimizing $\KL(d^+_{\pi,H}\Vert q)$ is maximizing $\Ent(d^+_{\pi,H})+\E_{d^+_{\pi,H}}[\log q]$: the entropy of the \emph{state occupancy} enters with a plus sign. It is not the entropy of the action distribution, and neither maximizes the other in general (\cref{ex:fpg-entropy}).
\end{proposition}
\begin{proof}
$\KL(d\Vert q)=\sum_{d>0}d\log d-\sum_{d>0}d\log q$, and every $\log q$ is finite on $\supp d$.
\end{proof}

\begin{proposition}[the fixed-target optimum; \cref{prop:fpg-opt}]\claim{FPG-8}{D}\label{prop:m-fpg-opt}
In a finite MDP with tabular softmax policies: $J\ge0$ with equality if and only if $d_\phi=q$; if $q=q_E$ and $J(\phi)=0$, the learner has the expert's $H$-step return for every state reward (\cref{prop:m-return}); if the learner can reach a state outside $\supp q_E$, $\KL(d_\phi\Vert q_E)=+\infty$ for every $\phi$; for $\JS$ and $\KL(q_E\Vert d_\phi)$, $\inf_\phi J=0$, but the infimum need not be attained with finite logits (\cref{ex:fpg-noattain}).
\end{proposition}

The third statement is why $\KL(d\Vert q)$ needs smoothing against a deterministic expert, which never visits most states.

\begin{remark}[FPG's Dirac-target theorems; \cref{rem:fpg-dirac}]\claim{FPG-C1}{C}\label{rem:m-dirac}
FPG's optimality results (Theorem~5.1, p.~6; Theorem~4.2, p.~5) use $D_f(p\Vert\delta_g)=f(p(g))+f'(\infty)(1-p(g))$, a monotone function of the goal visitation. For a non-degenerate expert target the divergence depends on the whole occupancy and the argument has no analogue. We draw no optimality guarantee from these theorems.
\end{remark}
